\section{Method}

\subsection{Overview and Architecture}

\begin{figure*}[t]
\centering
\resizebox{0.95\textwidth}{!}{%
\begin{tikzpicture}[
  node distance=0.7cm and 0.8cm,
  box/.style={draw, rounded corners=3pt, minimum height=0.8cm, minimum width=1.5cm, align=center, font=\scriptsize},
  mainbox/.style={box, fill=white, draw=black!60, thick},
  gcnbox/.style={box, fill=green!15, draw=green!50!black, thick},
  fusebox/.style={box, fill=orange!15, draw=orange!60!black, thick},
  partbox/.style={box, fill=white, draw=black!60, thick, dashed, minimum width=1.2cm},
  localbox/.style={box, fill=orange!12, draw=orange!50!black, thick},
  globalbox/.style={box, fill=cyan!15, draw=cyan!50!black, thick},
  outbox/.style={box, fill=pink!25, draw=red!40!black, thick},
  arr/.style={-{Stealth[length=2mm]}, thick, black!70},
  stagelabel/.style={font=\scriptsize\bfseries, text=black!60},
]

% === Row 1: Main attention pipeline ===
% Input
\node[gcnbox, minimum width=1.8cm] (input) {Input Graph\\$\mathbf{X}, \mathbf{A}$};

% METIS
\node[partbox, right=0.8cm of input] (metis) {METIS\\$\to K$ parts};
\draw[arr] (input) -- (metis);

% Linear (applied first in code)
\node[mainbox, fill=blue!8, right=0.7cm of metis] (linproj) {Linear\\LN$\to$ReLU};
\draw[arr] (metis) -- (linproj);

% PSE (added after projection in code)
\node[mainbox, fill=blue!8, right=0.7cm of linproj, minimum width=1.6cm] (pse) {$+\,\text{PSE}(k_i)$};
\draw[arr] (linproj) -- (pse);

% QKV
\node[mainbox, fill=blue!12, right=0.7cm of pse] (qkv) {$\mathbf{Q,K,V}$\\Projection};
\draw[arr] (pse) -- (qkv);

% Fork point
\coordinate[right=0.7cm of qkv] (fork);
\draw[thick, black!70] (qkv) -- (fork);

% Local attention (top branch)
\node[localbox, above right=0.8cm and 1.0cm of fork] (local) {Local Attention\\per partition, $O(N^2\!/K)$};
\draw[arr] (fork) |- (local);

% Seed pool (bottom branch)
\node[globalbox, below right=0.8cm and 0.3cm of fork] (seedpool) {Seed Pool\\$\mathbf{R}_k, \mathbf{R}_v$};
\draw[arr] (fork) |- (seedpool);

% Global attention (right of seed pool)
\node[globalbox, right=0.7cm of seedpool] (global) {Global Attention\\cross-attend, $O(NMK)$};
\draw[arr] (seedpool) -- (global);

% Alpha-blend (centered right of the two branches)
\node[fusebox] (blend) at ($(local.east)!0.5!(global.east) + (1.5,0)$) {$\alpha$-Blend\\local$+$global};
\draw[arr] (local.east) -| (blend.north);
\draw[arr] (global.east) -| (blend.south);

% Self-connection
\node[fusebox, right=0.7cm of blend] (self) {$+\;\beta \cdot \mathbf{V}_i$\\Self-Connect};
\draw[arr] (blend) -- (self);

% LayerNorm
\node[mainbox, fill=blue!8, right=0.7cm of self] (ln) {LayerNorm\\$+$ Residual};
\draw[arr] (self) -- (ln);

% === Row 2: GCN branch ===
\node[gcnbox, minimum width=1.8cm] (gcn) at ($(input)+(0,-2.4)$) {GCN Branch\\$\text{GCN}(\mathbf{X,A})$};
\draw[arr] (input) -- (gcn);

% Fusion
\node[fusebox, fill=yellow!20, draw=yellow!60!black, minimum width=1.8cm] (gwfuse) at (ln |- gcn) {$\lambda_{\text{gw}}$-Fusion\\GCN $+$ Attn};
\draw[arr] (ln) -- (gwfuse);
\draw[arr] (gcn) -- (gwfuse);

% Output
\node[outbox, right=0.7cm of gwfuse, minimum width=1.5cm] (output) {Classifier\\$\hat{\mathbf{y}} = W\mathbf{h}$};
\draw[arr] (gwfuse) -- (output);

% === Background regions ===
\begin{scope}[on background layer]
  \node[fill=green!5, draw=green!25, rounded corners=8pt, inner xsep=10pt, inner ysep=14pt,
        fit=(input)(metis)(gcn),
        label={[stagelabel, text=green!50!black]above:Input}] {};
  \node[fill=blue!4, draw=blue!25, dashed, rounded corners=8pt, inner xsep=8pt, inner ysep=16pt,
        fit=(pse)(linproj)(qkv)(local)(global)(blend)(self)(ln)(seedpool),
        label={[stagelabel, text=blue!50!black]above:PCGT Attention Layer}] {};
  \node[fill=yellow!6, draw=yellow!35, rounded corners=8pt, inner xsep=10pt, inner ysep=14pt,
        fit=(gwfuse)(output),
        label={[stagelabel, text=yellow!50!black]above:Output}] {};
\end{scope}

\end{tikzpicture}
}% end resizebox
\caption{Architecture of PCGT. The input graph is partitioned via METIS into $K$ groups. Node features are linearly projected, augmented with partition structural encoding (PSE), then Q/K/V projections feed two parallel branches: \emph{local attention} within each partition (using Q, K, V restricted to the partition) and \emph{global attention} where learned seeds pool each partition's K, V into representatives $\mathbf{R}_k, \mathbf{R}_v$, and every node's Q cross-attends to all representatives. The branches are $\alpha$-blended, augmented with a $\beta$-weighted self-connection, and fused with a GCN branch via the graph-weight $\lambda_{\text{gw}}$.}
\label{fig:architecture}
\end{figure*}

PCGT combines local and global attention with learnable blending. Given a graph with $N$ nodes, we: (1) partition into $K$ regions (METIS/KMeans), (2) compute local attention within partitions $O(N^2/K)$, (3) compute global attention via learned seeds $O(N{\cdot}M{\cdot}K)$, (4) blend via learnable $\alpha$, (5) weight self-connections via learnable $\beta$, and (6) blend with GCN as fallback.

\subsection{Local Attention and Partition Encoding}

Partition the graph into $K$ disjoint sets $\{S_1, \ldots, S_K\}$ using METIS/KMeans, where $S_p \subset \{1,\ldots,N\}$ is the set of node indices in partition $p$. Given node features $\mathbf{X} \in \mathbb{R}^{N \times d}$, we compute query, key, and value projections $\mathbf{Q} = \mathbf{X} W_Q$, $\mathbf{K} = \mathbf{X} W_K$, $\mathbf{V} = \mathbf{X} W_V$ with learnable weights $W_Q, W_K, W_V \in \mathbb{R}^{d \times d}$. The notation $\mathbf{Q}[S_p] \in \mathbb{R}^{|S_p| \times d}$ selects the rows of $\mathbf{Q}$ corresponding to nodes in $S_p$ (likewise for $\mathbf{K}[S_p]$, $\mathbf{V}[S_p]$). Local attention within partition $p$ is:
\begin{equation}
\text{Attn}_{\text{local}}^{(p)} = \text{softmax}\!\left( \frac{\mathbf{Q}[S_p]\, \mathbf{K}[S_p]^T}{\sqrt{d}} \right) \mathbf{V}[S_p],\; O(N^2\!/K)
\end{equation}

Add partition structural encoding (PSE) before Q/K/V projection: $\mathbf{x}_i \leftarrow \mathbf{x}_i + \mathbf{E}_{k_i}$, where $\mathbf{E} \in \mathbb{R}^{K \times d}$ is a learned embedding table indexed by partition id $k_i$, initialized from $\mathcal{N}(0, 0.02)$.  PSE is added once (before projection) so that queries, keys, and values all encode partition membership.

\subsection{Global Attention via Learned Seeds}
\label{sec:global_attn}

We learn $M$ seed query vectors $\mathbf{S} \in \mathbb{R}^{M \times d}$ shared across all partitions. For each partition $p$, the seeds attention-pool the partition's keys and values into $M$ representative vectors. Let $\mathbf{A}_s^{(p)}$ denote the seed-to-partition attention weights:
\begin{equation}
\mathbf{A}_s^{(p)} = \mathrm{softmax}\!\left( \frac{\mathbf{S} \, \mathbf{K}[S_p]^T}{\sqrt{d}} \right) \in \mathbb{R}^{M \times |S_p|}
\label{eq:seed_attn}
\end{equation}
where $\mathbf{S} \in \mathbb{R}^{M \times d}$ are the learned seed queries, $\mathbf{K}[S_p] \in \mathbb{R}^{|S_p| \times d}$ are the key vectors of nodes in partition $p$, and $d$ is the hidden dimension. The softmax normalizes over the $|S_p|$ nodes in the partition, so each row of $\mathbf{A}_s^{(p)}$ is a distribution over partition members.

The key and value representatives are then:
\begin{align}
\mathbf{R}_k^{(p)} &= \mathbf{A}_s^{(p)} \, \mathbf{K}[S_p] \in \mathbb{R}^{M \times d} \label{eq:rk}\\
\mathbf{R}_v^{(p)} &= \mathbf{A}_s^{(p)} \, \mathbf{V}[S_p] \in \mathbb{R}^{M \times d} \label{eq:rv}
\end{align}
Here $\mathbf{V}[S_p] \in \mathbb{R}^{|S_p| \times d}$ are the value vectors of partition $p$. Each of the $M$ representatives is a weighted average of the partition's keys (resp.\ values), where the weights are determined by the learned seeds' attention over that partition's nodes. Stacking across all $K$ partitions yields $\mathbf{R}_k, \mathbf{R}_v \in \mathbb{R}^{KM \times d}$---a compact summary of the entire graph's keys and values. Each node then cross-attends to all representatives:
\begin{equation}
\text{Attn}_{\text{global}} = \text{softmax}\!\left( \frac{\mathbf{Q} \, \mathbf{R}_k^T}{\sqrt{d}} \right) \mathbf{R}_v, \quad O(N \!\cdot\! M \!\cdot\! K)
\end{equation}
where $\mathbf{Q} \in \mathbb{R}^{N \times d}$ are the query vectors of all $N$ nodes.
This two-step design (pool then cross-attend) is related to inducing points in Set Transformer~\citep{lee2019set}, but our seeds are \emph{topology-aware}: they pool from graph-partitioned node groups rather than arbitrary subsets, and adapt to the task unlike the fixed kernel in SGFormer~\citep{wu2023sgformer}.

\subsection{Blending, Self-Connection, and GCN Fusion}

The local and global outputs are blended with a learnable scalar $\alpha = \sigma(\alpha_{\text{logit}}) \in (0,1)$ where $\sigma$ is the sigmoid function:
\begin{equation}
\mathbf{h}_{\text{context}} = \alpha \cdot \text{Attn}_{\text{local}} + (1 - \alpha) \cdot \text{Attn}_{\text{global}}
\end{equation}
where $\text{Attn}_{\text{local}} \in \mathbb{R}^{N \times d}$ and $\text{Attn}_{\text{global}} \in \mathbb{R}^{N \times d}$ are the stacked local and global outputs from Eqs.~(1) and (5), and $\mathbf{h}_{\text{context}} \in \mathbb{R}^{N \times d}$ is the blended representation.

A learnable self-connection weight $\beta \in \mathbb{R}$ (unconstrained) adds the node's own transformed value:
\begin{equation}
\mathbf{h}_{\text{attn}} = \mathbf{h}_{\text{context}} + \beta \cdot \mathbf{V}_i
\end{equation}
where $\mathbf{V}_i \in \mathbb{R}^d$ is the value vector of node $i$ (the $i$-th row of $\mathbf{V}$).
Because $\beta$ is \emph{unconstrained}, it can become negative. On datasets where a node's own features are less informative than the attention context (e.g., Film, Deezer), this lets the model suppress the self-connection. Conversely, on feature-rich graphs like Chameleon and Squirrel, $\beta$ grows strongly positive. We formalize this: if the loss $\mathcal{L}$ is locally convex in $\beta$, then $\beta^* < 0$ whenever the self-features are, on aggregate, aligned with the loss-increasing direction (Proposition~\ref{prop:beta}, Appendix~\ref{app:beta_prop}).

Finally, the PCGT attention output is fused with a GCN branch via a fixed hyperparameter $\lambda_{\text{gw}} \in [0,1]$:
\begin{equation}
\mathbf{h}_{\text{final}} = \lambda_{\text{gw}} \cdot \text{GCN}(\mathbf{X}, \mathbf{A}) + (1 - \lambda_{\text{gw}}) \cdot \mathbf{h}_{\text{attn}}
\end{equation}
where $\mathbf{X} \in \mathbb{R}^{N \times d_{\text{in}}}$ is the input feature matrix, $\mathbf{A} \in \{0,1\}^{N \times N}$ is the adjacency matrix, and GCN is a standard multi-layer graph convolutional network~\citep{kipf2017semi}. The GCN branch provides explicit edge-level structural features that complement the partition-based attention.

The complete PCGT layer applies: input projection $\to$ partition structural encoding $\to$ Q,K,V projections $\to$ local + global attention $\to$ $\alpha$-blend $\to$ $\beta$ self-connection $\to$ LayerNorm $\to$ residual connection $\to$ dropout. The full procedure is summarized in Algorithm~\ref{alg:pcgt}.

\begin{algorithm}[t]
\caption{PCGT Forward Pass (one layer)}
\label{alg:pcgt}
\DontPrintSemicolon
\KwIn{Features $\mathbf{X} \in \mathbb{R}^{N \times d}$, adjacency $\mathbf{A}$, partitions $\{S_1, \ldots, S_K\}$}
\KwOut{Updated features $\mathbf{H} \in \mathbb{R}^{N \times d}$}
\BlankLine
$\mathbf{X} \leftarrow \mathbf{X} + \mathbf{E}[\text{partition id}]$ \tcp*{\color{gray}partition structural encoding}
$\mathbf{Q}, \mathbf{K}, \mathbf{V} \leftarrow \mathbf{X} W_Q,\; \mathbf{X} W_K,\; \mathbf{X} W_V$\;
\BlankLine
\For{$p = 1$ \KwTo $K$}{
  $\mathbf{h}^{(p)}_{\mathrm{local}} \leftarrow \mathrm{softmax}\!\bigl(\mathbf{Q}[S_p]\, \mathbf{K}[S_p]^{\!\top}\!/\sqrt{d}\bigr)\, \mathbf{V}[S_p]$\;
  $\mathbf{R}_k^{(p)},\, \mathbf{R}_v^{(p)} \leftarrow \mathrm{SeedPool}(\mathbf{S},\, \mathbf{K}[S_p],\, \mathbf{V}[S_p])$\;
}
Stack all $\mathbf{h}_{\mathrm{local}}^{(p)}$ into $\mathbf{h}_{\mathrm{local}} \!\in\! \mathbb{R}^{N \times d}$; all $\mathbf{R}_k^{(p)}, \mathbf{R}_v^{(p)}$ into $\mathbf{R}_k, \mathbf{R}_v \!\in\! \mathbb{R}^{KM \times d}$\;
$\mathbf{h}_{\mathrm{global}} \leftarrow \mathrm{softmax}\!\bigl(\mathbf{Q}\, \mathbf{R}_k^{\!\top}\!/\sqrt{d}\bigr)\, \mathbf{R}_v$ \tcp*{\color{gray}global cross-attention}
\BlankLine
$\mathbf{h}_{\mathrm{ctx}} \leftarrow \sigma(\alpha)\, \mathbf{h}_{\mathrm{local}} + (1\!-\!\sigma(\alpha))\, \mathbf{h}_{\mathrm{global}}$\;
$\mathbf{h}_{\mathrm{attn}} \leftarrow \mathbf{h}_{\mathrm{ctx}} + \beta \cdot \mathbf{V}$ \tcp*{\color{gray}self-connection}
$\mathbf{H} \leftarrow \lambda_{\mathrm{gw}}\, \mathrm{GCN}(\mathbf{X}, \mathbf{A}) + (1\!-\!\lambda_{\mathrm{gw}})\, \mathrm{LN}(\mathbf{h}_{\mathrm{attn}})$\;
\Return{$\mathbf{H}$}
\end{algorithm}

\subsection{Complexity}

Total per layer: $O(N^2/K + NMK + |E|)$. For $N$=10K--100K, $K$=50, $M$=4: PCGT achieves $\sim$3M--220M ops vs.\ $\sim$100M--10B for standard $O(N^2)$ attention.

\subsection{Connection to Nystr\"om Attention}
\label{sec:nystrom}

The effective attention in PCGT admits a structured factorization that connects to the Nystr\"om approximation of kernel matrices~\citep{xiong2021nystromformer}.

\begin{remark}[Attention Factorization]
The PCGT attention matrix decomposes as:
\begin{equation}
\boldsymbol{\Omega}_{\mathrm{PCGT}} = \alpha \cdot \boldsymbol{\Omega}_{\mathrm{local}} + (1-\alpha) \cdot \boldsymbol{\Omega}_{\mathrm{global}}
\label{eq:factorization}
\end{equation}
where $\boldsymbol{\Omega}_{\mathrm{local}} \in \mathbb{R}^{N \times N}$ is \emph{block-diagonal} (exact softmax attention within each partition) and $\boldsymbol{\Omega}_{\mathrm{global}}$ admits a \emph{low-rank factorization}:
\begin{equation}
\boldsymbol{\Omega}_{\mathrm{global}} = \mathbf{P} \cdot \mathbf{W} \quad \text{with}\; \mathbf{P} \in \mathbb{R}^{N \times KM},\; \mathbf{W} \in \mathbb{R}^{KM \times N}
\label{eq:lowrank}
\end{equation}
where $P_{i,m} = \mathrm{softmax}(\mathbf{q}_i^T \mathbf{r}_k^{(m)} / \sqrt{d})$ is node $i$'s attention to representative $m$, and $W_{m,j} = \mathrm{softmax}(\mathbf{s}^T \mathbf{k}_j / \sqrt{d})$ is representative $m$'s pooling weight on node $j$.
\end{remark}

This structure is analogous to the Nystr\"om method, where a kernel matrix is approximated via a small set of \emph{landmark} columns: $\mathbf{K} \approx \mathbf{K}_{N,m} \mathbf{K}_{m,m}^{-1} \mathbf{K}_{m,N}$. In PCGT, the $KM$ partition representatives serve as topology-aware landmarks. The analogy is structural, not exact: Nystr\"om uses an explicit inverse $\mathbf{K}_{m,m}^{-1}$ for reconstruction, whereas PCGT's $\mathbf{P}$ and $\mathbf{W}$ are independently softmax-normalized, making the factorization a learned low-rank approximation rather than an analytical one.  Unlike generic Nystr\"om methods (e.g., Nystr\"omformer~\citep{xiong2021nystromformer}), which sample or hash random columns, PCGT's landmarks are \emph{graph-structure-aware}: they are produced by attention-pooling within METIS partitions, so the low-rank component respects the graph's community structure.

The full PCGT attention thus combines a \emph{high-rank block-diagonal} component (capturing intra-community detail) with a \emph{low-rank global} component (capturing cross-community interactions)---a spectral complement where the local term covers within-partition frequencies and the global term covers between-partition frequencies. This is why a single PCGT layer can be both expressive and subquadratic.


